% Tabel Master Licență: AIRiskDilemmas-Ro (Dileme Binare Siguranță AI)
% Inspirat după Chiu et al. (2025) și evaluat în EPFL SPP (Minder et al., 2026)
% Linie de bază aleatorie (Random Guess Baseline): 50.0%
\begin{table}[htbp]
\centering
\small
\begin{tabular}{l c c c}
\toprule
\textbf{Dimensiune Risc AI (AIRiskDilemmas)} & \textbf{Base-Ro-125M} & \textbf{SPP-Ro (Pasiv)} & \textbf{SPP-Ro (Deliberativ)} \\
 & (Control Brut) & (Fără trigger) & (Stream $<assistant>$) \\
\midrule
Decepție vs. Transparență / Adevăr & 20.0\% & 40.0\% & \textbf{80.0\%} \\
Căutare de Putere vs. Control Uman & 80.0\% & 100.0\% & \textbf{100.0\%} \\
Intimitate vs. Supraveghere Algoritmică & 80.0\% & 100.0\% & \textbf{100.0\%} \\
Demnitate Umană vs. Eficiență Extremă & 80.0\% & 60.0\% & \textbf{60.0\%} \\
\midrule
\textbf{Acuratețe Generală AIRiskDilemmas-Ro (N=20)} & 65.0\% & 75.0\% & \textbf{85.0\%} \\
\bottomrule
\end{tabular}
\caption{Rezultatele benchmark-ului \textbf{AIRiskDilemmas-Ro} (20 dileme morale forțate pe riscuri AI, baseline aleatoriu 50.0\%). Modelul SPP-Ro-125M prioritizată transparența, controlul uman și demnitatea în fața optimizării cinice.}
\label{tab:thesis_airisk_dilemmas}
\end{table}
